\section*{Chapitre \ref{ch:b1:topology} --- Topologie de la droite réelle}

\begin{solution}{exo:b1:topology:1}
$\intoo{0}{1} \cup \intoo{2}{3}$~: \omterm{def:b1:topology:open}{ouvert} (réunion d'\omterm{def:b1:topology:open}{ouverts}), non
\omterm{def:b1:topology:closed}{fermé} ($\frac 1n \to 0$, qui est en dehors).

$\intco{0}{1}$~: ni l'un ni l'autre. Non \omterm{def:b1:topology:open}{ouvert} (aucun \omterm{prop:b1:reals:intervals}{intervalle}
autour de $0$ n'y est contenu)~; non \omterm{def:b1:topology:closed}{fermé} ($1 - \frac1n \to 1 \notin$
l'\omterm{def:b1:logic:sets}{ensemble}).

$\{0\} \cup \intcc{1}{2}$~: \omterm{def:b1:topology:closed}{fermé} (réunion finie de \omterm{def:b1:topology:closed}{fermés}), non
\omterm{def:b1:topology:open}{ouvert} (cela échoue en $0$).

$\R \setminus \Z$~: \omterm{def:b1:topology:open}{ouvert} ($\Z$ est \omterm{def:b1:topology:closed}{fermé}), non \omterm{def:b1:topology:closed}{fermé}~: la suite
$\bigl(\frac 1n\bigr)$ y est à valeurs, mais sa limite $0$ appartient
à $\Z$, c'est-à-dire s'échappe de l'\omterm{def:b1:logic:sets}{ensemble}.

$\Q \cap \intoo{0}{1}$~: ni l'un ni l'autre. Non \omterm{def:b1:topology:open}{ouvert}~: tout
\omterm{prop:b1:reals:intervals}{intervalle} autour d'un rationnel contient des irrationnels. Non
\omterm{def:b1:topology:closed}{fermé}~: il contient des suites tendant vers l'irrationnel
$\frac{\sqrt 2}{2}$ (\omterm{def:b1:topology:dense}{densité}).
\end{solution}

\begin{solution}{exo:b1:topology:2}
$A = \intoc{0}{1} \cup \{2\}$~: $\mathring A = \intoo{0}{1}$,
$\overline A = \intcc{0}{1} \cup \{2\}$, $\partial A = \{0, 1, 2\}$.

$A = \R \setminus \Q$~: $\mathring A = \emptyset$ (tout \omterm{prop:b1:reals:intervals}{intervalle}
contient des rationnels), $\overline A = \R$ (\omterm{def:b1:topology:dense}{densité} des
irrationnels), $\partial A = \R$.

$A = \bigl\{\frac{(-1)^n n}{n+1}\bigr\}$~: les termes pairs tendent
vers $1$, les termes impairs vers $-1$, et aucune de ces deux valeurs
n'appartient à $A$. $\mathring A = \emptyset$ (points isolés),
$\overline A = A \cup \{-1, 1\}$, $\partial A = \overline A$.
\end{solution}

\begin{solution}{exo:b1:topology:3}
\emph{Par le complémentaire.} $F = \{a_1 < a_2 < \dots < a_k\}$~: le
complémentaire est la réunion des \omterm{prop:b1:reals:intervals}{intervalles} \omterm{def:b1:topology:open}{ouverts}
$\intoo{-\infty}{a_1}$, $\intoo{a_i}{a_{i+1}}$, $\intoo{a_k}{+\infty}$
--- \omterm{def:b1:topology:open}{ouvert} d'après la \cref{prop:b1:topology:openstable}.

\emph{Par les suites.} Soit $u_n \in F$, $u_n \to \ell$. Avec
$\varepsilon = \min\{\abs{a_i - a_j} : i \neq j\}/2 > 0$ (ou
n'importe quel $\varepsilon$ si $F$ est un singleton)~: à partir d'un
certain rang, tous les termes sont à distance $\leq \varepsilon$ de
$\ell$, donc à distance $\leq 2\varepsilon$ les uns des autres, ce qui
les force à être un seul et même $a_i$ à partir de ce rang~; alors
$\ell = a_i \in F$.
\end{solution}

\begin{solution}{exo:b1:topology:4}
$\subseteq$~: $\overline A \cup \overline B$ est \omterm{def:b1:topology:closed}{fermé} (réunion finie)
et contient $A \cup B$, donc il contient le \emph{plus petit} \omterm{def:b1:topology:closed}{fermé}
contenant $A \cup B$, à savoir $\overline{A \cup B}$. $\supseteq$~:
$A \subseteq A \cup B$ donne $\overline A \subseteq \overline{A \cup B}$
(l'\omterm{def:b1:topology:closure}{adhérence} est croissante~: les points adhérents à $A$ le sont au
plus grand \omterm{def:b1:logic:sets}{ensemble}), et de même pour $B$.

Contre-exemple pour les intersections~: $A = \intoo{0}{1}$, $B =
\intoo{1}{2}$~: $\overline{A \cap B} = \overline\emptyset = \emptyset$
mais $\overline A \cap \overline B = \{1\}$.
\end{solution}

\begin{solution}{exo:b1:topology:5}
Utilisons la caractérisation séquentielle
(\cref{thm:b1:topology:seqclosed}). Posons $S = \{u_n\} \cup \{\ell\}$
et soit $(v_k)$ une suite de $S$ avec $v_k \to m$~; montrons que
$m \in S$. Deux cas. Si une valeur $v \in S$ est prise une infinité de
fois par $(v_k)$, une \omterm{def:b1:seq:subsequence}{suite extraite} constante donne $m = v \in S$.
Sinon, chaque valeur n'est prise qu'un nombre fini de fois~; en
particulier, pour chaque $n$, le terme $u_n$ n'apparaît qu'un nombre
fini de fois, et $\ell$ aussi. Alors, pour tout $N$, les indices $k$
tels que $v_k \in \{u_0, \dots, u_N, \ell\}$ sont en nombre fini~: les
$v_k$ restants sont des termes $u_n$ avec $n > N$. Étant donné
$\varepsilon > 0$, choisissons $N$ tel que $\abs{u_n - \ell} \leq
\varepsilon$ pour $n > N$~: tous les $v_k$ sauf un nombre fini
vérifient $\abs{v_k - \ell} \leq \varepsilon$. Donc $v_k \to \ell$, et
$m = \ell \in S$.
\end{solution}

\begin{solution}{exo:b1:topology:6}
$U + A = \bigcup_{a \in A} (U + a)$, et chaque translaté $U + a$ est
\omterm{def:b1:topology:open}{ouvert} (il suffit de translater les certificats d'\omterm{prop:b1:reals:intervals}{intervalles}). Une
réunion d'\omterm{def:b1:topology:open}{ouverts} est \omterm{def:b1:topology:open}{ouverte} (\cref{prop:b1:topology:openstable}).

\omterm{def:b1:topology:closed}{Fermés}~: $\Z$ et $\sqrt 2\,\Z$ sont \omterm{def:b1:topology:closed}{fermés} (comme tout $\alpha\Z$~:
les suites convergentes y sont stationnaires, cf.\
\cref{ex:b1:topology:closed}). Leur somme $\Z + \sqrt 2\,\Z$ est \omterm{def:b1:topology:dense}{dense}
dans $\R$ (\cref{exo:b1:reals:9}) mais n'est pas $\R$ (elle est
dénombrable, ou plus simplement~: $\frac{\sqrt 2}{2} \notin \Z +
\sqrt2\Z$, sans quoi $\sqrt 2$ serait rationnel --- écrire
$\frac{\sqrt2}{2} = m + n\sqrt 2$ force $(2n - 1)\sqrt 2 = -2m$, donc
$\sqrt 2 \in \Q$ sauf si $n = \frac12$, ce qui est impossible). Une
partie \omterm{def:b1:topology:dense}{dense} et propre n'est pas \omterm{def:b1:topology:closed}{fermée}~: son \omterm{def:b1:topology:closure}{adhérence} est $\R$,
distincte d'elle-même.
\end{solution}

\begin{solution}{exo:b1:topology:7}
$\Z$~: le \omterm{def:b1:topology:open}{voisinage} $\intoo{n - \frac12}{n + \frac12}$ de $n$
rencontre $\Z$ en $n$ seulement.

$A = \{\frac1n : n \in \N^*\}$~: autour de $\frac 1n$, l'\omterm{prop:b1:reals:intervals}{intervalle} de
rayon $\frac{1}{n} - \frac{1}{n+1} = \frac{1}{n(n+1)}$ (divisé par
deux, disons) le sépare de ses voisins --- tous les points sont
isolés. Pourtant $0 \in \overline A \setminus A$~: des points isolés
n'empêchent pas des adhérents venus de l'extérieur.

Si tous les points de $A$ sont isolés~: aucun point de $A$ n'est
\omterm{def:b1:topology:closure}{intérieur}, car un point \omterm{def:b1:topology:closure}{intérieur} possède autour de lui tout un
\omterm{prop:b1:reals:intervals}{intervalle} de points de $A$ (un \omterm{prop:b1:reals:intervals}{intervalle} est infini), ce qui
contredit l'isolement. Donc $\mathring A = \emptyset$.
\end{solution}

\begin{solution}{exo:b1:topology:8}
Si $A \subseteq \intcc{-M}{M}$, le \omterm{def:b1:topology:closed}{fermé} $\intcc{-M}{M}$ contient $A$,
donc contient $\overline A$ (plus petit \omterm{def:b1:topology:closed}{fermé} contenant $A$)~:
$\overline A$ est bornée.

Soit $s = \sup A$ (fini). Comme $A \subseteq \overline A$, on a $\sup
\overline A \geq s$. Réciproquement, $\overline A \subseteq
\intoc{-\infty}{s}$~: cette demi-droite est \omterm{def:b1:topology:closed}{fermée} et contient $A$~;
donc tout élément de $\overline A$ est $\leq s$, ce qui donne
$\sup\overline A \leq s$. D'où l'égalité.

$s \in \overline A$~: par la caractérisation en $\varepsilon$
(\cref{prop:b1:reals:epsilon}), tout \omterm{prop:b1:reals:intervals}{intervalle} $\intoo{s -
\varepsilon}{s + \varepsilon}$ contient un élément de $A$~: $s$ est
adhérent.
\end{solution}

\begin{solution}{exo:b1:topology:9}
Supposons $A$ \omterm{def:b1:topology:open}{ouvert} et \omterm{def:b1:topology:closed}{fermé}, avec $a \in A$ et $b \in \R
\setminus A$~; quitte à échanger les rôles, $a < b$. L'\omterm{def:b1:logic:sets}{ensemble} $B = A
\cap \intcc{a}{b}$ est non vide (il contient $a$) et borné~: posons
$s = \sup B$. D'après l'\cref{exo:b1:topology:8}, $s \in \overline B
\subseteq \overline A = A$ ($A$ \omterm{def:b1:topology:closed}{fermé}). Remarquons que $s \leq b$ et,
puisque $b \notin A$~: $s < b$. Or $A$ est \omterm{def:b1:topology:open}{ouvert}~: un \omterm{prop:b1:reals:intervals}{intervalle}
$\intoo{s - r}{s + r}$ est contenu dans $A$, et on peut prendre
$r < b - s$. Alors $s + \frac{r}{2}$ appartient à $A \cap
\intcc{a}{b} = B$ et dépasse $s$ --- ce qui contredit $s = \sup B$. Il
n'existe donc aucun tel couple $(a, b)$~: l'un des \omterm{def:b1:logic:sets}{ensembles} $A$,
$\R \setminus A$ est vide.
\end{solution}

\begin{solution}{exo:b1:topology:10}
$I_x$ est une réunion d'\omterm{prop:b1:reals:intervals}{intervalles} \omterm{def:b1:topology:open}{ouverts} contenant tous $x$~: il
est \omterm{def:b1:topology:open}{ouvert}, et c'est un \omterm{prop:b1:reals:intervals}{intervalle}, car il est convexe --- si
$u < z < v$ avec $u, v \in I_x$, alors $u$ et $v$ appartiennent à des
\omterm{prop:b1:reals:intervals}{sous-intervalles} \omterm{def:b1:topology:open}{ouverts} $J_u \ni x$, $J_v \ni x$ de $U$, et
$J_u \cup J_v$ est un \omterm{prop:b1:reals:intervals}{intervalle} (tous deux contiennent $x$) inclus
dans $U$ et contenant $z$~; donc $z \in I_x$
(\cref{prop:b1:reals:intervals}).

Si $I_x \cap I_y \neq \emptyset$~: $I_x \cup I_y$ est alors un
\omterm{prop:b1:reals:intervals}{intervalle} \omterm{def:b1:topology:open}{ouvert} (convexe~: deux \omterm{prop:b1:reals:intervals}{intervalles} qui se chevauchent)
contenu dans $U$ et contenant $x$ et $y$, donc $I_x \cup I_y \subseteq
I_x$ et $\subseteq I_y$ par maximalité de chacun~: $I_x = I_y$.

Ainsi $U$ est la réunion disjointe des \omterm{def:b1:logic:sets}{ensembles} $I_x$ distincts
(chaque $x \in U$ appartient à son propre $I_x$). Dénombrabilité~:
chaque \omterm{prop:b1:reals:intervals}{intervalle} \omterm{def:b1:topology:open}{ouvert} non vide $I$ de la famille contient un
rationnel $q_I$ (\cref{thm:b1:reals:density}), et des \omterm{prop:b1:reals:intervals}{intervalles}
distincts, donc disjoints, reçoivent des rationnels distincts~: la
famille s'injecte dans $\Q$, qui est dénombrable (il est indexé par
des couples d'entiers). Il y a donc au plus une infinité dénombrable
d'\omterm{prop:b1:reals:intervals}{intervalles}.
\end{solution}

\begin{solution}{exo:b1:topology:11}
$\overline A = A \cup A'$. ($\supseteq$) On a toujours $A \subseteq
\overline A$~; et si $x \in A'$, tout \omterm{def:b1:topology:open}{voisinage} de $x$ rencontre $A
\setminus \{x\} \subseteq A$, donc $x$ est adhérent. ($\subseteq$)
Soit $x \in \overline A$. Si $x \in A$, c'est fini. Si $x \notin A$,
tout \omterm{def:b1:topology:open}{voisinage} de $x$ rencontre $A = A \setminus \{x\}$~: $x \in
A'$.

Par conséquent, $A$ est \omterm{def:b1:topology:closed}{fermé} $\iff$ $A = \overline A = A \cup A'$
$\iff$ $A' \subseteq A$.

$A = \{\frac 1n\}$~: $0$ est point d'accumulation ($\frac 1n \to
0$, termes $\neq 0$)~; chaque $\frac 1n$ est isolé
(\cref{exo:b1:topology:7}), donc n'est pas dans $A'$~; et un point
$x \notin A \cup \{0\}$ possède tout un \omterm{prop:b1:reals:intervals}{intervalle} évitant $A$
(entre les deux voisins de $x$ dans $A \cup \{0\}$, ou au-delà de
$1$). Donc $A' = \{0\}$.

$\Z' = \emptyset$~: tout entier est isolé, tout non-entier possède
un \omterm{def:b1:topology:open}{voisinage} contenu dans $\R \setminus \Z$.

$\Q' = \R$~: tout \omterm{prop:b1:reals:intervals}{intervalle} autour d'un réel quelconque contient
une infinité de rationnels (\cref{thm:b1:reals:density}), en
particulier un différent du centre.
\end{solution}

\begin{solution}{exo:b1:topology:12}
\begin{enumerate}
  \item Pour tout $a \in A$~: $\abs{x - a} \leq \abs{x - y} +
        \abs{y - a}$, donc $d_A(x) \leq \abs{x - y} + \abs{y - a}$~;
        en passant à la \omterm{def:b1:reals:bounds}{borne inférieure} sur $a$~: $d_A(x) \leq
        \abs{x - y} + d_A(y)$. En échangeant $x$ et $y$ on obtient
        l'autre inégalité~: $\abs{d_A(x) - d_A(y)} \leq \abs{x - y}$.
  \item $d_A(x) = 0$ $\iff$ pour tout $\varepsilon > 0$ il existe
        $a \in A$ avec $\abs{x - a} < \varepsilon$ $\iff$ tout
        \omterm{prop:b1:reals:intervals}{intervalle} autour de $x$ rencontre $A$ $\iff$ $x \in \overline A$.
        Si $F$ est \omterm{def:b1:topology:closed}{fermé} et $x \notin F = \overline F$, alors
        $d_F(x) \neq 0$, c'est-à-dire $d_F(x) > 0$.
  \item Si $d_F(x) < \varepsilon$, posons $r = \varepsilon -
        d_F(x) > 0$~: pour $\abs{y - x} < r$, le point (1) donne
        $d_F(y) \leq d_F(x) + \abs{x - y} < \varepsilon$~:
        l'\omterm{prop:b1:reals:intervals}{intervalle} $\intoo{x - r}{x + r}$ est contenu dans
        $V_\varepsilon$, qui est donc \omterm{def:b1:topology:open}{ouvert}~; il contient $F$
        puisque $d_F = 0$ sur $F$. Enfin $x \in
        \bigcap_{\varepsilon>0} V_\varepsilon$ $\iff$ $d_F(x) <
        \varepsilon$ pour tout $\varepsilon$ $\iff$ $d_F(x) = 0$
        $\iff$ $x \in \overline F = F$. Comme
        $\bigcap_{\varepsilon > 0} V_\varepsilon =
        \bigcap_{n \geq 1} V_{1/n}$, tout \omterm{def:b1:topology:closed}{fermé} est une
        intersection dénombrable d'\omterm{def:b1:topology:open}{ouverts}.
\end{enumerate}
\end{solution}

\begin{solution}{pb:b1:topology:1}
\textbf{1.} $C_1 = \intcc{0}{\frac13} \cup \intcc{\frac23}{1}$
et
\[
C_2 = \intcc{0}{\tfrac19} \cup \intcc{\tfrac29}{\tfrac13}
\cup \intcc{\tfrac23}{\tfrac79} \cup \intcc{\tfrac89}{1} .
\]
Récurrence~: si $C_n$ est une réunion disjointe de $2^n$ segments
\omterm{def:b1:topology:closed}{fermés} de longueur $3^{-n}$, supprimer le tiers médian \omterm{def:b1:topology:open}{ouvert} de
chacun laisse deux segments \omterm{def:b1:topology:closed}{fermés} de longueur $3^{-n-1}$ par
parent~: $2^{n+1}$ segments, deux à deux disjoints (les enfants de
parents distincts sont séparés parce que les parents l'étaient~;
les enfants d'un même parent sont séparés par le trou retiré).

\textbf{2.} Chaque $C_n$ est une réunion finie de segments, donc
\omterm{def:b1:topology:closed}{fermé}~; $C = \bigcap C_n$ est une intersection de \omterm{def:b1:topology:closed}{fermés}~: c'est un
\omterm{def:b1:topology:closed}{fermé} (\cref{def:b1:topology:closed})~; borné ($\subseteq
\intcc{0}{1}$)~: donc compact d'après le \cref{thm:b1:topology:compact}.
Non vide~: $0$ appartient au segment le plus à gauche de chaque
$C_n$. Soit $a$ une extrémité d'un segment $S$ de $C_n$. Pour
$m \leq n$, $a \in C_n \subseteq C_m$. Pour les étapes suivantes~: la
suppression du tiers médian n'enlève jamais une extrémité, et $a$
est de nouveau une extrémité de l'un des deux enfants de $S$ (celui
qui touche $a$)~; par récurrence $a \in C_m$ pour tout $m \geq n$~:
$a \in C$.

\textbf{3.} Longueur totale de $C_n$~: $2^n \cdot 3^{-n} =
(\frac23)^n \to 0$. Étant donné $\varepsilon > 0$, choisissons $n$
tel que $(\frac23)^n \leq \varepsilon$~: alors $C \subseteq C_n$,
réunion d'un nombre fini de segments de longueur totale $\leq
\varepsilon$.

\textbf{4.} Soit $I \subseteq C$ un \omterm{prop:b1:reals:intervals}{intervalle} contenant deux points
distincts. Pour tout $n$~: $I \subseteq C_n$, et $I$, étant convexe,
doit être contenu dans un \emph{seul} segment de $C_n$ (rencontrer
deux segments forcerait $I$ à contenir un point du trou qui les
sépare, lequel est hors de $C_n$). Donc la longueur de $I$ est $\leq
3^{-n}$ pour tout $n$~: contradiction. Ainsi les seuls \omterm{prop:b1:reals:intervals}{intervalles}
contenus dans $C$ sont vides ou réduits à un point~; en particulier
aucun $\intoo{x-r}{x+r}$ ne tient dans $C$~: $\mathring C =
\emptyset$. Comme $C$ est \omterm{def:b1:topology:closed}{fermé}, $\overline C = C$ est d'\omterm{def:b1:topology:closure}{intérieur}
vide~: $C$ est nulle part \omterm{def:b1:topology:dense}{dense}.

\textbf{5.} Posons $\varphi_0(x) = \frac x3$ et $\varphi_2(x) =
\frac{2 + x}{3}$, \omterm{def:b1:logic:inj}{bijections} affines croissantes de $\intcc{0}{1}$
sur $\intcc{0}{\frac13}$ et $\intcc{\frac23}{1}$. Affirmation~:
$C_{n+1} = \varphi_0(C_n) \cup \varphi_2(C_n)$. Pour $n = 0$ c'est
la question 1. Récurrence~: une \omterm{def:b1:logic:map}{application} affine croissante envoie
le tiers médian d'un segment sur le tiers médian du segment image,
donc la suppression des tiers médians commute avec $\varphi_0$ et
$\varphi_2$~; en appliquant l'étape de suppression à $C_{n+1} =
\varphi_0(C_n) \cup \varphi_2(C_n)$ on obtient $C_{n+2} =
\varphi_0(C_{n+1}) \cup \varphi_2(C_{n+1})$. En intersectant sur
$n$~: pour $x \leq \frac13$, $x \in C \iff x \in \varphi_0(C_n)$
pour tout $n$ $\iff 3x \in \bigcap C_n = C$~; de même sur
$\intcc{\frac23}{1}$~; et aucun point de $\intoo{\frac13}{\frac23}$
n'appartient à $C_1$. Donc $C = \varphi_0(C) \cup \varphi_2(C)$, de
façon disjointe.

\textbf{6.} Récurrence sur $n$~; le cas $n = 0$ dit que tout $x \in
\intcc{0}{1}$ possède un code, ce qui est le \cref{pb:b1:reals:1}
(question 9 pour $x < 1$~; $1 = (0.\overline 2)_3$). Supposons
l'équivalence au rang $n$. Si $x \in C_{n+1}$~: d'après la question
5, $x = \varphi_i(z)$ avec $z \in C_n$ et $i \in \{0, 2\}$~; si
$(e_k)$ est un code de $z$ dont les $n$ premiers chiffres sont sans
$1$, alors $(i, e_1, e_2, \dots)$ a pour sommes partielles $\frac i3
+ \frac13\sum_{k\leq m} e_k 3^{-k} \to \varphi_i(z) = x$~: c'est un
code de $x$ dont les $n + 1$ premiers chiffres sont sans $1$.
Réciproquement, si $x$ possède un code $(d_k)$ avec $d_1, \dots,
d_{n+1} \in \{0, 2\}$~: la suite décalée $(d_2, d_3, \dots)$ a une
certaine valeur $z \in \intcc{0}{1}$, ses $n$ premiers chiffres sont
sans $1$, et le calcul des sommes partielles lu à l'envers donne
$x = \varphi_{d_1}(z)$~; par hypothèse de récurrence $z \in C_n$,
donc $x \in C_{n+1}$ par la question 5. Enfin~: un code entièrement
sans $1$ place $x$ dans chaque $C_n$, donc dans $C$~; réciproquement
si $x \in C$, alors pour chaque $n$ l'un des au plus deux codes de
$x$ (\cref{pb:b1:reals:1}, question 11) a ses $n$ premiers chiffres
sans $1$~; un même code fixé doit convenir pour des $n$
arbitrairement grands (\omterm{cor:b1:counting:pigeonhole}{principe des tiroirs} entre deux codes), et un
code dont les $n$ premiers chiffres sont sans $1$ pour des $n$
arbitrairement grands est sans $1$ tout court.

\textbf{7.} $0 = (0.\overline 0)_3$, $1 = (0.\overline 2)_3$,
$\frac13 = (0.0\overline{2})_3$ (le jumeau impropre de $(0.1)_3$),
$\frac23 = (0.2\overline{0})_3$. Sommes géométriques~:
\[
(0.\overline{02})_3 = \sum_{j\geq1} \frac{2}{9^{\,j}}
= \frac{2/9}{1 - 1/9} = \frac14 ,
\qquad
(0.\overline{20})_3 = \sum_{j\geq1} \frac{2}{3\cdot 9^{\,j-1}}
= \frac{2/3}{1 - 1/9} = \frac34 ,
\]
tous deux sans $1$~: $\frac14, \frac34 \in C$. (Les sommes infinies
abrègent des \omterm{def:b1:reals:bounds}{bornes supérieures} de sommes partielles, comme dans
le \cref{pb:b1:reals:1}.) Les extrémités des segments de $C_n$ sont de
la forme $m/3^n$ (récurrence~: les extrémités des enfants sont des
extrémités du parent, ou en diffèrent d'un multiple de
$3^{-n-1}$). Si $\frac14 = \frac{m}{3^n}$ alors $3^n = 4m$, et
$4 \nmid 3^n$~: impossible. Donc $\frac14 \in C$ sans jamais être
une extrémité.

\textbf{8.} Supposons que $x$ ait deux codes sans $1$ distincts. Avoir
deux codes signifie (\cref{pb:b1:reals:1}, question 11, en base
$3$) que $x = m/3^N \in \intoo{0}{1}$ et que les deux codes sont~:
le code fini, dont le dernier chiffre non nul $d_N \in \{1, 2\}$ est
suivi de $0$, et son jumeau, portant $d_N - 1$ à la position $N$
puis des $2$. Si $d_N = 1$ le premier contient un $1$~; si $d_N
= 2$ le jumeau porte $d_N - 1 = 1$. Dans les deux cas, au plus l'un
des deux est sans $1$~: contradiction. Donc chaque $x \in C$ possède
exactement un code sans $1$ (l'existence vient de la question 6), et
des suites $\{0,2\}$ distinctes ont des valeurs distinctes. Toute
suite $\{0,2\}$ a une valeur dans $\intcc{0}{1}$ (sommes partielles
$\leq 1$) dont tous les préfixes sont sans $1$, donc une valeur dans
chaque $C_n$, c'est-à-dire dans $C$~: l'\omterm{def:b1:logic:map}{application} «~valeur~» est
une \omterm{def:b1:logic:inj}{bijection} de l'\omterm{def:b1:logic:sets}{ensemble} des suites $\{0,2\}$ sur $C$.

\textbf{9.} Soit $(d^{(k)})$ le code sans $1$ de $x_k$ et posons
$e_k = 2 - d^{(k)}_k \in \{0, 2\}$~: c'est une suite $\{0,2\}$ dont
la valeur $y$ appartient à $C$ et admet $(e_k)$ pour unique code sans
$1$ (question 8). Pour chaque $k$, les codes de $y$ et de $x_k$
diffèrent à la position $k$, donc $y \neq x_k$~: aucune \omterm{def:b1:logic:map}{application}
$\N^* \to C$ n'est \omterm{def:b1:logic:inj}{surjective}. Un \omterm{def:b1:logic:sets}{ensemble} non dénombrable de
longueur nulle~: énorme en \omterm{def:b1:counting:card}{cardinal}, minuscule en mesure ---
simultanément.

\textbf{10.} Compacité~: le caractère \omterm{def:b1:topology:closed}{fermé} de l'intersection
infinie, plus le caractère borné (question 2) --- la seule propriété
qui ne soit pas portée par une simple inclusion. Longueur nulle~:
$C \subseteq C_n$, de longueur totale $(\frac23)^n$ (question 3).
Nulle part \omterm{def:b1:topology:dense}{dense}~: $C \subseteq C_n$ force les \omterm{prop:b1:reals:intervals}{intervalles} contenus
dans $C$ à être de longueur $\leq 3^{-n}$ (question 4). La
non-dénombrabilité ne repose sur aucune inclusion~: elle réclame
toute la structure de l'intersection, encodée dans la \omterm{def:b1:logic:inj}{bijection} de
la question 8.

\textbf{11.} Remplaçons $d_n$ par $2 - d_n$~: la nouvelle suite est
encore une suite $\{0,2\}$, donc sa valeur $x_n$ appartient à $C$~;
les sommes partielles au-delà du rang $n$ diffèrent exactement de
$2\cdot3^{-n}$, donc $\abs{x_n - x} = 2\cdot3^{-n}$. Ainsi $x_n \neq
x$ et $x_n \to x$~: tout point de $C$ est limite d'autres points de
$C$ --- $C$ est \emph{parfait}, sans point isolé.

\textbf{12.} D'après la récurrence de la question 5, les segments de
$C_n$ sont exactement les $\intcc{t}{t + 3^{-n}}$ où $t$ parcourt
les valeurs des suites $\{0,2\}$ de longueur $n$. Étant donné $x \in
C$ de code $(d_k)$, la troncature $t_n$ (chiffres $d_1 \dots d_n$
puis des $0$) est donc une extrémité gauche, et $0 \leq x - t_n
\leq 3^{-n}$~: les extrémités sont \omterm{def:b1:topology:dense}{denses} dans $C$. Elles forment
une partie de $\{m/3^n : m, n\}$, \omterm{def:b1:logic:sets}{ensemble} indexé par des couples
d'entiers, donc dénombrable (comme pour $\Q$ dans
l'\cref{exo:b1:topology:10}). Comme $C$ n'est pas dénombrable (question
9), tous les points de $C$ sauf une infinité dénombrable ne sont pas
des extrémités --- $\frac14$ (question 7) est la partie émergée de
cet iceberg.

\textbf{13.} Choisissons $n$ tel que $3^{-n} < y - x$. Les deux
points $x, y$ sont dans $C_n$, et ils ne peuvent pas appartenir au
même segment (de longueur $3^{-n} < y - x$)~: le trou retiré entre
leurs segments fournit un $z$ avec $x < z < y$ et $z \notin C_n
\supseteq C$. Ainsi deux points quelconques de $C$ sont séparés par
le complémentaire~: les seules parties convexes de $C$ sont les
singletons --- $C$ est totalement discontinu.

\textbf{14.} Si $(d_k)$ est le code sans $1$ de $x$, la suite
$(2 - d_k)$ est encore une suite $\{0,2\}$, de sommes partielles
\[
\sum_{k=1}^{n} (2 - d_k)3^{-k} = (1 - 3^{-n}) - \sum_{k=1}^n
d_k 3^{-k} \longrightarrow 1 - x :
\]
donc $1 - x \in C$. Ainsi $1 - C \subseteq C$, et en appliquant
l'\omterm{def:b1:logic:map}{application} deux fois on obtient $1 - C = C$~: l'\omterm{pb:b1:topology:1}{ensemble de
Cantor} est symétrique par rapport à $\frac12$.

\textbf{15.} Les sommes partielles vérifient $t_n \to x$ et
$t'_n \to x'$ (une suite croissante converge vers sa \omterm{def:b1:reals:bounds}{borne
supérieure}, c'est-à-dire vers la valeur), donc $t_n + t'_n \to x + x'$
d'après le \cref{thm:b1:seq:operations}. Étant donné $y \in
\intcc{0}{1}$ de code $(e_k)$, choisissons $(a_k, b_k) = (0,0),
(0,2), (2,2)$ selon que $e_k = 0, 1, 2$~: alors $a_k + b_k = 2e_k$,
les suites $(a_k)$, $(b_k)$ sont des suites $\{0,2\}$ de valeurs $x,
x' \in C$, et
\[
x + x' = \lim_n\,(t_n + t'_n) = \lim_n 2\sum_{k=1}^n e_k 3^{-k}
= 2y :
\]
tout $y \in \intcc{0}{1}$ est le milieu de deux points de $C$.

\textbf{16.} La question 15 donne $\intcc{0}{2} = 2\,\intcc{0}{1}
\subseteq C + C$, et $C + C \subseteq \intcc{0}{1} +
\intcc{0}{1} = \intcc{0}{2}$~: d'où l'égalité. Puis, en utilisant
$1 - C = C$~:
\[
C - C = C + (C - 1) = (C + C) - 1 = \intcc{-1}{1} .
\]
Cas concret~: $1 = \frac14 + \frac34$, somme de deux éléments de $C$
qui ne sont pas des extrémités.

\textbf{17.} La longueur mesure la portion de droite qu'occupe
l'\omterm{def:b1:logic:sets}{ensemble} lui-même~; elle ne dit rien de l'\omterm{def:b1:logic:sets}{ensemble} des
\emph{sommes}, qui est l'image de la famille à deux paramètres
$C \times C$ par $(x, x') \mapsto x + x'$ --- les deux suites de
chiffres sont choisies indépendamment, et c'est exactement cette
liberté qui remplit $\intcc{0}{2}$. Aucun théorème ne majore la
longueur d'un \omterm{def:b1:logic:sets}{ensemble} somme par les longueurs des termes, et $C$ est
la preuve qu'aucun ne le peut.

\textbf{18.} D'après le \cref{pb:b1:reals:1} (question 18), $x$ est
rationnel si et seulement si son développement propre est périodique
à partir d'un certain rang. Le code sans $1$ de $x \in C$ est soit ce
développement propre, soit le jumeau impropre d'un développement
fini~; or un développement fini et son jumeau (des $2$ à partir d'un
rang) sont tous deux périodiques à partir d'un certain rang, de sorte
que la périodicité du code sans $1$ équivaut à la rationalité de $x$.
Division de $\frac1{13}$ en base $3$ ($r_0 = 1$)~: $3 = 13\cdot0 +
3$, $9 = 13\cdot0 + 9$, $27 = 13\cdot2 + 1$, et le reste revient à
$1$~: chiffres $\overline{002}$, donc $\frac1{13} =
(0.\overline{002})_3$, sans $1$ et périodique~: un élément rationnel
de $C$. (Vérification~: $\frac{2/27}{1 - 1/27} = \frac{2}{26} =
\frac1{13}$.)

\textbf{19.} La suite valant $d_k = 2$ aux positions triangulaires
$k = \frac{j(j+1)}{2}$ et $0$ ailleurs est une suite $\{0,2\}$, donc
sa valeur $x^*$ appartient à $C$ (question 8). Elle comporte une
infinité de $2$ séparés par des écarts $j + 1 \to \infty$ entre deux
consécutifs, donc elle n'est pas périodique à partir d'un certain
rang (une période $T$ finirait par imposer des $2$ à des écarts
$\leq T$~: l'argument des écarts croissants du \cref{pb:b1:reals:1},
question 20)~; d'après la question 18, $x^* \notin \Q$. Et d'après la
question 9 jointe à la dénombrabilité de $\Q$, tous les éléments de
$C$ sauf une infinité dénombrable sont irrationnels~: $x^*$ est la
norme, non l'exception.

\textbf{20.} Soit $y \in \intcc{0}{1}$~: il possède un code binaire
$(c_k)$ avec $c_k \in \{0, 1\}$ (\cref{pb:b1:reals:1}, question
9, en base $2$~; $y = 1$ prend la suite constante égale à $1$). Alors
$(2c_k)$ est une suite $\{0,2\}$, sa valeur $x$ appartient à $C$, et
$h(x)$ est la valeur de $(c_k)$, à savoir $y$~: $h$ envoie $C$ sur
$\intcc{0}{1}$ tout entier. Si $C$ était l'image d'une \omterm{def:b1:logic:map}{application}
définie sur $\N^*$, la composition avec $h$ énumérerait tout
$\intco{0}{1}$, ce qui contredirait le théorème diagonal du
\cref{pb:b1:reals:1} (question 22)~: $C$ n'est pas dénombrable, de
nouveau. Un \omterm{def:b1:logic:sets}{ensemble} de longueur nulle qui se surjecte sur un segment
tout entier.

\textbf{21.} D'après la question 5, $C$ est la réunion disjointe de
$\varphi_0(C)$ et $\varphi_2(C)$, chacun étant un translaté de la
copie réduite $\frac13 C$. L'additivité, l'homothétie et l'invariance
par translation donnent
\[
L(C) = L(\varphi_0(C)) + L(\varphi_2(C))
= \tfrac13 L(C) + \tfrac13 L(C) = \tfrac23\,L(C),
\]
donc $\frac13 L(C) = 0$~: $L(C) = 0$. L'autosimilarité seule condamne
$C$ à la longueur nulle --- la question 3 n'a fait qu'exécuter la
sentence.

\textbf{22.} Un segment de longueur $l_n$ perd un \omterm{prop:b1:reals:intervals}{intervalle} central
de longueur $4^{-(n+1)}$, ce qui laisse deux segments de longueur
$l_{n+1} = \frac{l_n - 4^{-(n+1)}}{2}$~; partant de $l_0 = 1$, une
récurrence confirme $l_n = \frac{2^n + 1}{2\cdot4^n}$~: en effet
$\frac12\Bigl(\frac{2^n+1}{2\cdot4^n} - \frac{1}{4^{n+1}}\Bigr)
= \frac{2(2^n + 1) - 1}{2\cdot4^{n+1}} = \frac{2^{n+1} +
1}{2\cdot4^{n+1}}$, et $l_n > 0$ toujours~: la construction ne
s'épuise jamais. $K = \bigcap K_n$ est \omterm{def:b1:topology:closed}{fermé} et borné, donc compact~;
un \omterm{prop:b1:reals:intervals}{intervalle} contenu dans $K$ est contenu dans un segment de $K_n$,
de longueur $l_n \to 0$~: \omterm{def:b1:topology:closure}{intérieur} vide. Longueur retirée~:
$\sum_{n\geq0} 2^n \cdot 4^{-(n+1)} =
\frac14\sum_{n\geq0}\bigl(\frac12\bigr)^n = \frac12$, et chaque $K_n$
a pour longueur totale $2^n l_n = \frac{2^n + 1}{2^{n+1}} >
\frac12$. Soit maintenant un nombre fini d'\omterm{prop:b1:reals:intervals}{intervalles} \omterm{def:b1:topology:open}{ouverts} de
réunion $U \supseteq K$. D'après l'énoncé compagnon de
l'\cref{ex:b1:topology:nested}, $U \supseteq K_n$ pour un certain $n$~;
en admettant l'additivité de la longueur sur les réunions finies
d'\omterm{prop:b1:reals:intervals}{intervalles}, la longueur totale des \omterm{prop:b1:reals:intervals}{intervalles} du recouvrement est
au moins celle de $K_n$, qui dépasse $\frac12$. Ainsi $K$ est nulle
part \omterm{def:b1:topology:dense}{dense}, et pourtant aucun recouvrement bon marché n'existe~:
petitesse topologique (nulle part \omterm{def:b1:topology:dense}{dense}) et petitesse métrique
(longueur nulle) sont deux notions réellement différentes, et $K$ les
sépare.

\textbf{23.} Atteinte~: posons $d = d(x, F)$ et choisissons $a_k \in
F$ avec $\abs{x - a_k} \leq d + \frac1k$~: les $a_k$ sont bornés,
donc Bolzano--Weierstrass (\cref{thm:b1:seq:bw}) en extrait
$a_{\varphi(k)} \to a$, avec $a \in F$ ($F$ \omterm{def:b1:topology:closed}{fermé},
\cref{thm:b1:topology:seqclosed}) et $\abs{x - a} = \lim
\abs{x - a_{\varphi(k)}} = d$. Passons au maximum~: si $y \in C$,
$d(y, C) = 0$~; sinon $y$ appartient à un trou retiré à une étape
$n \geq 1$, \omterm{prop:b1:reals:intervals}{intervalle} \omterm{def:b1:topology:open}{ouvert} de longueur $3^{-n}$ dont les deux
extrémités appartiennent à $C$ (question 2), donc $d(y, C) \leq
\frac{3^{-n}}{2} \leq \frac16$, avec égalité seulement si $n = 1$ et
$y$ au centre du trou $\intoo{\frac13}{\frac23}$, c'est-à-dire
$y = \frac12$~; et en effet $d\bigl(\frac12, C\bigr) = \frac16$
puisque $C \cap \intoo{\frac13}{\frac23} = \emptyset$ et
$\frac13, \frac23 \in C$. Donc $\max_{y\in\intcc{0}{1}} d(y, C)
= \frac16$, atteint exactement en $\frac12$.

\textbf{24.} Les extrémités forment une partie dénombrable \omterm{def:b1:topology:dense}{dense} de
$C$ (question 12)~: énumérons-les en une seule suite
$(e_j)_{j\geq1}$, à valeurs dans $C$. Ses valeurs d'\omterm{def:b1:topology:closure}{adhérence}
appartiennent toutes à $C$ ($C$ \omterm{def:b1:topology:closed}{fermé}). Réciproquement, fixons $x \in
C$~: pour chaque $n$, les segments des $C_m$ contenant $x$ ($m \geq
n$) ont leurs extrémités à distance $\leq 3^{-m} \leq 3^{-n}$ de $x$,
donc une infinité d'extrémités distinctes sont à distance $\leq
3^{-n}$ de $x$~; choisissons des indices $j_1 < j_2 < \dots$ avec
$\abs{e_{j_n} - x} \leq 3^{-n}$~: voilà une \omterm{def:b1:seq:subsequence}{suite extraite}
convergeant vers $x$. Ainsi l'\omterm{def:b1:logic:sets}{ensemble} des valeurs d'\omterm{def:b1:topology:closure}{adhérence} de
$(e_j)$ est exactement $C$ --- une seule suite s'accumulant en une
infinité non dénombrable de points, l'extrême opposé d'une suite
convergente, dont l'\omterm{def:b1:logic:sets}{ensemble} des valeurs d'\omterm{def:b1:topology:closure}{adhérence} est un
singleton.

\textbf{25.} (i) La construction a consommé~: la stabilité des \omterm{def:b1:topology:closed}{fermés}
par intersection quelconque (existence de $C$ comme \omterm{def:b1:topology:closed}{fermé}), le
théorème de compacité \cref{thm:b1:topology:compact} (questions 2,
22, 23), et les caractérisations séquentielles du caractère \omterm{def:b1:topology:closed}{fermé} et
de l'\omterm{def:b1:topology:closure}{adhérence} (l'argument des compacts emboîtés de
l'\cref{ex:b1:topology:nested} et la question 23). (ii) Les quatre
couples~: longueur nulle et pourtant non dénombrable (questions 3,
9)~; \omterm{def:b1:topology:closed}{fermé} et pourtant d'\omterm{def:b1:topology:closure}{intérieur} vide (question 4)~; parfait ---
sans point isolé --- et pourtant totalement discontinu (questions 11,
13)~; négligeable et pourtant $C + C = \intcc{0}{2}$ (question 16).
(iii) La \omterm{def:b1:logic:inj}{surjection} $h$ de la question 20, rendue continue et
croissante, devient l'escalier du diable dans la théorie des
fonctions continues~; et dans la théorie de la mesure du volume de
Licence 3, $C$ est le témoin standard que «~négligeable~» ne signifie
pas «~dénombrable~», son cousin gras (question 22) séparant «~nulle
part \omterm{def:b1:topology:dense}{dense}~» de «~négligeable~».
\end{solution}
